\documentclass[11pt]{article}
\usepackage[margin=1in]{geometry}
\usepackage[hidelinks]{hyperref}
\pagestyle{empty}
\setlength{\parskip}{6pt}
\setlength{\parindent}{0pt}

\begin{document}

To: Prof. Xizhao Wang, Editor-in-Chief, \textit{International Journal of
Machine Learning and Cybernetics} (Springer)

Dear Prof. Wang,

We are pleased to submit our manuscript, ``FreqLite: A Lightweight
Frequency-Decomposed Linear Model with Adaptive Reversible Normalization for
Robust Long-Term Time-Series Forecasting,'' for consideration as an original
research article in \textit{International Journal of Machine Learning and
Cybernetics}.

Long-term time-series forecasting is a core task in engineering applications
spanning energy systems, traffic management, weather monitoring, and industrial
process control. Our work directly targets a practical gap: existing accurate
forecasters (Transformer-based) are too resource-intensive for constrained
deployment environments, while cheap linear models have limited representational
capacity. FreqLite bridges this gap with a principled frequency-domain
architecture that is both highly accurate and exceptionally resource-efficient.

The manuscript makes three contributions. First, \textbf{FreqLite}, an
ultra-lightweight, channel-independent frequency-decomposed linear forecaster: a
learnable partition-of-unity spectral filter splits the input signal into
frequency bands modeled by dedicated linear heads. Evaluated across seven
datasets (including high-variate Electricity with 321 channels and Traffic with
862 channels) against eight baselines, FreqLite at long lookback ($L{=}336$)
attains the best overall average test MSE and the most per-setting wins of any
model, \textbf{outperforming iTransformer (ICLR 2024) by 4.3\% and TimesNet
(ICLR 2023) by 23.2\%} in paired per-cell MSE---while using $13\times$ less
peak GPU memory and $18\times$ less training time than TimesNet---all on a
single 4\,GB consumer laptop GPU. These improvements are statistically
significant under paired Wilcoxon tests ($p < 10^{-4}$). Second, \textbf{Adaptive Reversible
Instance Normalization (A-RevIN)}, a regime-adaptive reversible normalization
that strictly generalizes RevIN, engaging under non-stationarity (validated on
ILI and a controlled distribution-shift stress test) and reducing to vanilla
RevIN without harm on stationary data. Third, a \textbf{comprehensive,
reproducible accuracy--efficiency study} on commodity hardware across 7 datasets
and 8 baselines, with ablations isolating each component and a distribution-shift
stress test that directly probes A-RevIN's robustness claims.

This work is aligned with IJMLC's focus on the hybrid development of machine
learning methodology and its practical application: all experiments run on a
single 4\,GB laptop GPU, the complete codebase is publicly available and
reproducible, and the paper reports limitations and non-wins honestly.

This manuscript is original, has not been published previously, and is not under
consideration elsewhere. All authors have approved the submission; the authors
declare no competing interests and received no specific funding.

Thank you for your consideration.

\vspace{6pt}
Sincerely,

\vspace{6pt}
Mirza Samad Ahmed Baig (Corresponding author) \\
Fandaqah, Al Khobar, Saudi Arabia \\
Mirza@fandaqah.com

\vspace{4pt}
Syeda Anshrah Gillani \\
Heidelberg University, Germany \\
syeda.gilani@stud.uni-heidelberg.de

\vspace{4pt}
Asher Ali \\
Fandaqah, Al Khobar, Saudi Arabia \\
CTO@Fandaqah.com

\end{document}
